\documentclass[a4paper,11pt]{article}
\usepackage[a4paper,margin=2.5cm]{geometry}
\usepackage[T1]{fontenc}
\usepackage[indonesian]{babel}
\usepackage{booktabs}
\usepackage{amsmath}

\title{Tahap 3: Conceptual Modeling\\
\large Rumusan Model Konseptual DSS CCTV Jatinangor}
\date{}
\begin{document}
\maketitle

\noindent\textit{Status: DRAF perumusan (diisi nilai lapangan oleh tim; merujuk sintesis Tahap 2 v2, protokol 2021--2026).}

\section{Alternatif (level segmen gang, mengikuti \#39 dan \#31)}

\begin{table}[htbp]
\centering
\small
\setlength{\tabcolsep}{3pt}
\begin{tabular}{p{1.0cm}p{3.2cm}p{1.8cm}p{2.8cm}p{5.0cm}}
\toprule
\textbf{Kode} & \textbf{Segmen (contoh operasional, verifikasi lapangan oleh tim)} & \textbf{Jangkar wilayah} & \textbf{Ciri segmen} & \textbf{Dasar pemilihan} \\
\midrule
A1 & Gang kos kawasan Sayang & Desa Sayang & Padat, akses sempit & Analog hunian padat \#126; logika coverage titik \#39 \\
A2 & Ruas jalan Cikeruh dekat gerbang kampus & Desa Cikeruh & Koridor utama kampus-kos & Volume mobilitas \#31; sampling titik kritis \#46 \\
A3 & Persimpangan Hegarmanah & Desa Hegarmanah & Simpul multi-cabang & Derajat adjacency \#39; peta risiko spasial \#31 \\
A4 & Gang Cipacing & Desa Cipacing & Minim PJU & Infrastruktur/penerangan \#83; intervensi terukur \#11 \\
A5 & Ruas Cileles & Desa Cileles & Jalur mobilitas malam & Flow/delay \#46; guna lahan sekitar kampus \#31 \\
\bottomrule
\end{tabular}
\caption{Alternatif segmen gang A1--A5}
\end{table}

Aturan: 4--6 alternatif; tiap alternatif satu segmen jalan/gang yang dapat dipasangi 1 unit CCTV; nama di atas contoh operasional, nama final dari observasi lapangan.

\section{Matriks kriteria C1--C9 (arah: skor tinggi = makin prioritas)}

Seluruh kriteria berarah benefit terhadap urgensi (skor tinggi = makin prioritas dipasangi CCTV).

\begin{table}[htbp]
\centering
\small
\setlength{\tabcolsep}{3pt}
\begin{tabular}{p{0.9cm}p{2.3cm}p{3.8cm}p{1.4cm}p{3.6cm}p{2.0cm}}
\toprule
\textbf{Kode} & \textbf{Kriteria} & \textbf{Definisi operasional} & \textbf{Satuan} & \textbf{Sumber data Jatinangor (utama / fallback)} & \textbf{Rujukan literatur} \\
\midrule
C1 & Riwayat kerawanan segmen & Kejadian kriminalitas/kecelakaan per km per tahun & kejadian/km/thn & Polsek atau laporan warga / observasi + wawancara RT & \#83, \#14, \#96, \#31 \\
C2 & Volume mobilitas malam & Orang lewat pukul 18.00--24.00 & orang/jam & Counting manual / CCTV eksisting (\#46: sampling titik kritis) & \#31, \#46, \#96 \\
C3 & Defisiensi penerangan & Tingkat kekurangan cahaya PJU (skor 1--5 survei) & skor 1--5 & Survei malam / data Dishub (fallback: observasi) & \#83, \#11 \\
C4 & Efisiensi coverage & Estimasi luas ter-cover per titik CCTV & m2/titik & Pemetaan + uji coverage model \#39 & \#39, \#56 \\
C5 & Kepadatan penghuni & Kamar kos per km gang & kamar/km & Pendataan kos / BPS desa & \#56, \#31, \#126 \\
C6 & Derajat simpul jalan & Jumlah cabang persimpangan & cabang & OpenStreetMap + observasi (mengikuti \#39) & \#39 \\
C7 & Coverage gap CCTV eksisting & Jarak ke CCTV terdekat & meter & Survei titik CCTV eksisting & \#39, \#56 \\
C8 & Kedekatan guna lahan rentan & Kos + ATM + warung 24 jam + gerbang kampus radius 150 m & unit & Observasi + POI (mengikuti logika \#31, \#83) & \#31, \#83, \#37 \\
C9 & Kelayakan implementasi & Kesiapan tiang/listrik/internet dikurangi estimasi biaya (skor) & skor & Survei PLN/jaringan + RAB (mengikuti \#14 cost, \#39 efisiensi) & \#14, \#39, \#83 \\
\bottomrule
\end{tabular}
\caption{Matriks kriteria keselamatan C1--C9}
\end{table}

\subsection{Frekuensi dukungan kriteria (disalin dari sintesis Tahap 2, Bagian B.3)}

\begin{table}[htbp]
\centering
\begin{tabular}{p{5cm}p{5cm}p{3cm}}
\toprule
\textbf{Kriteria} & \textbf{Muncul di paper} & \textbf{Frekuensi} \\
\midrule
C1 Riwayat kerawanan & \#83, \#14, \#96, \#31 & 4 paper \\
C2 Volume mobilitas malam & \#31, \#46, \#96 & 3 paper \\
C5 Kepadatan penghuni & \#56, \#31, \#126 & 3 paper \\
C8 Kedekatan guna lahan rentan & \#31, \#83, \#37 & 3 paper \\
C9 Kelayakan implementasi & \#14, \#39, \#83 & 3 paper \\
C3 Defisiensi penerangan & \#83, \#11 & 2 paper \\
C4 Efisiensi coverage & \#39, \#56 & 2 paper \\
C7 Coverage gap CCTV eksisting & \#39, \#56 & 2 paper \\
C6 Derajat simpul jalan & \#39 & 1 paper (sumber tunggal, dicatat apa adanya) \\
\bottomrule
\end{tabular}
\caption{Frekuensi dukungan kriteria}
\end{table}

\subsection{Urutan prioritas kriteria dan dasar penempatan}

Urutan ditentukan berdasarkan frekuensi kemunculan di 12 paper acuan; kriteria yang muncul di lebih banyak paper ditempatkan lebih tinggi. Ketika frekuensi sama, seri diputus berdasarkan relevansi langsung terhadap keputusan penempatan CCTV (kendala anggaran dan paparan didahulukan atas faktor pendukung).

\begin{table}[htbp]
\centering
\begin{tabular}{p{1.2cm}p{3.5cm}p{9.5cm}}
\toprule
\textbf{Urutan} & \textbf{Kriteria} & \textbf{Dasar penempatan} \\
\midrule
1 & C1 Riwayat kerawanan & Frekuensi tertinggi (4 paper); inti urgensi pemasangan di seluruh literatur safety \\
2 & C9 Kelayakan implementasi & Jangkar budget constraint \#39 (tiang 62 ke 33); tanpa kelayakan, prioritas tak dapat dieksekusi \\
3 & C2 Volume mobilitas malam & Paparan utama segmen (volume \#31, flow \#46); menentukan siapa yang dilindungi \\
4 & C5 Kepadatan penghuni & Populasi terpapar per segmen; legitimasi hunian padat \#126 \\
5 & C8 Kedekatan guna lahan rentan & Kedekatan target kriminalitas (kos, ATM, warung 24 jam); logika guna lahan \#31 \\
6 & C4 Efisiensi coverage & Luas lindung per unit alat; efisiensi \#39 dan trade-off coverage \#56 \\
7 & C7 Coverage gap CCTV eksisting & Menghindari duplikasi; blind-zone \#39 dan backup coverage \#56 \\
8 & C3 Defisiensi penerangan & Faktor pendukung; CCTV tetap butuh cahaya minimum (\#83, \#11) \\
9 & C6 Derajat simpul jalan & Bersumber tunggal \#39; tetap dipakai sebagai ciri persimpangan, dicatat apa adanya \\
\bottomrule
\end{tabular}
\caption{Urutan prioritas kriteria C1--C9}
\end{table}

\section{Skema pembobotan AHP multistakeholder (mengikuti \#37, acuan \#83 dan \#74)}

\subsection{Langkah perhitungan (per kelompok responden)}

\begin{enumerate}
\item Responden tiga kelompok: (a) pemerintah/Dinas (Dishub, Satpol PP), (b) kampus/pengelola kawasan, (c) warga (RT, pemilik kos).
\item Kuesioner pairwise comparison 9 kriteria per kelompok dengan skala Saaty 1--9 (1 = sama penting, 9 = mutlak lebih penting; nilai kebalikan untuk arah sebaliknya).
\item Bentuk matriks perbandingan \(A\) (\(n \times n\), \(n = 9\)), normalisasi per kolom, lalu vektor prioritas \(w\) = rata-rata baris matriks ternormalisasi.
\item Uji konsistensi: hitung vektor terbobot \(Aw\), lalu lamda-maks = rata-rata \((Aw)_i / w_i\); \(CI = (\text{lamda-maks} - n)/(n-1)\); \(CR = CI/RI\) dengan RI untuk \(n = 9\) sebesar 1,45. Syarat \(CR < 0{,}1\) (acuan \#83: CR 0,046).
\item Agregasi antar kelompok menggunakan geometric mean per sel matriks, lalu ulangi langkah 3--4 pada matriks agregat.
\item Validasi silang opsional: bobot objektif Entropy Weight Method dari data lapangan (mengikuti \#56) sebagai pembanding terhadap bobot subjektif AHP; selisih besar memicu diskusi ulang dengan pakar, bukan penggantian otomatis.
\end{enumerate}

\subsection{Contoh hitung ILUSTRASI (bukan data lapangan)}

Tiga kriteria (C1, C2, C9) dengan penilaian ilustrasi: C1 vs C2 = 3; C1 vs C9 = 2; C9 vs C2 = 2. Matriks \(A\):

\begin{table}[htbp]
\centering
\begin{tabular}{c|ccc}
\toprule
 & C1 & C2 & C9 \\
\midrule
C1 & 1 & 3 & 2 \\
C2 & 1/3 & 1 & 1/2 \\
C9 & 1/2 & 2 & 1 \\
\bottomrule
\end{tabular}
\caption{Matriks pairwise ilustrasi}
\end{table}

Jumlah kolom: C1 = 1,833; C2 = 6; C9 = 3,5. Normalisasi lalu rata-rata baris menghasilkan bobot \(w = (0{,}539;\ 0{,}164;\ 0{,}297)\); jumlah = 1,000.

Uji konsistensi: \(Aw = (1{,}625;\ 0{,}492;\ 0{,}894)\); \((Aw)_i/w_i = (3{,}015;\ 3{,}004;\ 3{,}009)\); lamda-maks = 3,009; \(CI = (3{,}009-3)/2 = 0{,}0046\); RI (\(n = 3\)) = 0,58; \(CR = 0{,}0046/0{,}58 = 0{,}008 < 0{,}1\) sehingga konsisten.

Contoh ini hanya menunjukkan cara hitung; matriks \(9 \times 9\) sesungguhnya diisi dari kuesioner lapangan oleh tim.

\section{Perankingan TOPSIS (mengikuti pola pemenang \#14)}

\subsection{Langkah perhitungan}

\begin{enumerate}
\item Matriks keputusan \(X\) (A1--A5 terhadap C1--C9) dari data lapangan; seluruh kriteria berarah benefit terhadap urgensi.
\item Normalisasi: \(r_{ij} = x_{ij}/\sqrt{\sum_i x_{ij}^2}\) per kolom.
\item Matriks terbobot: \(v_{ij} = w_j \times r_{ij}\) dengan \(w\) dari AHP.
\item Solusi ideal positif \(A^+\) (maksimum per kolom) dan negatif \(A^-\) (minimum per kolom).
\item Jarak tiap alternatif ke \(A^+\) (\(D^+\)) dan ke \(A^-\) (\(D^-\)), lalu closeness \(C_i = D^-_i/(D^+_i + D^-_i)\); peringkat dari \(C\) terbesar.
\item Analisis sensitivitas: ubah tiap bobot \(\pm\)20\% satu per satu (mengikuti \#83), hitung ulang ranking; peringkat dinyatakan robust bila tidak berubah posisi.
\end{enumerate}

\subsection{Contoh hitung ILUSTRASI (data dummy, bukan data lapangan)}

Tiga segmen (S1, S2, S3) dinilai 1--5 pada C1, C2, C9 dengan bobot ilustrasi \(w = (0{,}539;\ 0{,}164;\ 0{,}297)\) dari Seksi 3.2. Matriks dummy \(X\): S1 = (4, 3, 5); S2 = (2, 4, 3); S3 = (5, 2, 4).

Normalisasi: penyebut kolom = (6,708; 5,385; 7,071). Matriks terbobot \(V\): S1 = (0,321; 0,091; 0,210); S2 = (0,161; 0,122; 0,126); S3 = (0,402; 0,061; 0,168).

\(A^+ = (0{,}402;\ 0{,}122;\ 0{,}210)\); \(A^- = (0{,}161;\ 0{,}061;\ 0{,}126)\). Jarak: \(D^+ = (0{,}086;\ 0{,}255;\ 0{,}074)\); \(D^- = (0{,}184;\ 0{,}061;\ 0{,}245)\). Closeness: S1 = 0,682; S2 = 0,192; S3 = 0,768. Ranking: S3 \textgreater{} S1 \textgreater{} S2 (S3 unggul pada C1 dan C9 yang berbobot besar; S2 terendah meski C2-nya tertinggi, menunjukkan kompromi antar kriteria bekerja).

\section{Arsitektur konseptual (pola GIS-MCDM, merujuk \#31, \#11, \#46)}

\begin{verbatim}
INPUT (C1-C9: laporan warga, counting, survei PJU/kos/CCTV, OSM, RAB)
  -> ENGINE: [GIS (pemetaan segmen + gap)] + [AHP (bobot)] + [TOPSIS (ranking)]
  -> OUTPUT: ranking prioritas A1-A5 + peta prioritas gang + rekomendasi alokasi bertahap sesuai budget
\end{verbatim}

\begin{table}[htbp]
\centering
\begin{tabular}{p{2.5cm}p{4.5cm}p{4cm}p{3.5cm}}
\toprule
\textbf{Lapisan} & \textbf{Fungsi} & \textbf{Masukan} & \textbf{Keluaran} \\
\midrule
Input & Kumpulkan data diskrit per segmen & Rekap insiden, counting, survei, OSM, RAB & Dataset mentah C1--C9 \\
Preprocessing & Bersihkan dan kodifikasi & Dataset mentah & Matriks A1--A5 x C1--C9 \\
Decision Engine & Bobot + ranking + peta & Matriks + penilaian pakar & Bobot, ranking, peta prioritas \\
Output & Rekomendasi alokasi & Ranking + peta + budget & Urutan pemasangan bertahap \\
\bottomrule
\end{tabular}
\caption{Lapisan arsitektur konseptual}
\end{table}

\section{Yang ditolak dan ditunda (argumen tercatat)}

\begin{itemize}
\item Ditolak: simulasi operasional ala \#45 (butuh kalibrasi mahal); BWM-CoCoSo ala \#83 (selisih kinerja kecil, pakar tak tersedia).
\item Ditunda: response time/egress ala \#45 (data tak tersedia) menjadi keterbatasan + riset lanjutan.
\item Dicatat: template AHP+GIS+Fuzzy TOPSIS \#121 relevan namun gugur aturan tahun 2021--2026 (bukti disiplin protokol).
\end{itemize}

\section{Justifikasi pemilihan metode}

\subsection*{Justifikasi AHP (pembobotan)}

AHP dipilih karena:
\begin{enumerate}
\item Kesesuaian jenis data: AHP bekerja dengan penilaian pakar, sesuai kondisi Jatinangor yang kaya pengetahuan lokal (aparat, pengelola kampus, warga) tetapi minim data sensor. Paper \#37 membuktikan AHP berjalan lintas kelompok stakeholder; \#74 membuktikan AHP berjalan dari data sekunder/dokumen.
\item Dukungan literatur: AHP dipakai primer di 3 dari 12 paper (\#37, \#14, \#74) dan menjadi pembobotan tersering di korpus.
\item Transparan dan teraudit: matriks pairwise, vektor bobot, dan CR terdokumentasi sehingga tiap angka dapat ditelusuri dan diuji ulang.
\item Tidak memerlukan infrastruktur mahal: kuesioner dan lembar kerja cukup, berbeda dengan simulasi \#45 yang butuh kalibrasi.
\end{enumerate}

\subsection*{Justifikasi TOPSIS (perankingan)}

TOPSIS dipilih karena:
\begin{enumerate}
\item Kesesuaian jenis data: TOPSIS bekerja dengan data objektif terukur hasil observasi segmen (counting, survei, pemetaan), tanpa memerlukan penilaian pakar tambahan setelah bobot ditetapkan.
\item Dukungan literatur: pola AHP (bobot) + TOPSIS (ranking) terbukti pada \#14; komparasi \#83 menempatkan AHP-TOPSIS dekat dengan metode kompleks (agreement 88,6\% vs 92,8\%).
\item Konsep intuitif: alternatif terbaik adalah yang terdekat dengan ideal positif dan terjauh dari ideal negatif, sehingga mudah dijelaskan ke pemangku kepentingan non-teknis.
\item Menangani kompromi benefit secara langsung: closeness coefficient menggabungkan seluruh kriteria dalam satu skor peringkat, cocok untuk alokasi bertahap sesuai anggaran.
\end{enumerate}

\section{Pemetaan ke paper IEEE}

\begin{itemize}
\item Seksi IV-A: alternatif (Seksi 1) + matriks C1--C9 + frekuensi dan urutan prioritas (Seksi 2).
\item Seksi IV-B: skema AHP + TOPSIS + contoh hitung (Seksi 3--4) + justifikasi metode (Seksi 7).
\item Seksi IV-C: diagram arsitektur + tabel lapisan (Seksi 5) + justifikasi adopsi/penolakan (Seksi 6, merujuk B.1--B.4 Tahap 2).
\end{itemize}

\end{document}
